%!TEX encoding = UTF-8 Unicode
%!TEX root = ../lect-w12.tex

%%%

\ifkompendium\else
\begin{SlideExtra}{Valfri fördjupning: Jämförelse Scala och Java}
\begin{itemize}
  \item Föreläsningsserien är avslutad: läsveckorna w12-w14 inga föreläsningar.
  \item Appendix J: Introduktion till Java 
  \begin{itemize}
    \item Valfri jämförelse Scala--Java: Teori, Övning, Lab
    \item Föreberedelse inför efterföljande kurs: Programmering fördjupningskurs (pfk)
  \end{itemize}
\end{itemize}
\end{SlideExtra}
\fi 

% \ifkompendium\else

% \Subsection{Nästa vecka: repetition}

% \begin{SlideExtra}{Vad är grumligt? Vad är du nyfiken på?}\SlideFontSmall

% \begin{itemize}
% \item Skriv önskemål till de sista föreläsningarna på Discord i kanalen \texttt{\#önska} 
% \item Skriv ämnen/koncept ur kursen som du önskar att jag förklarar mer om -- sådant som fortfarande känns \Alert{grumligt}. 
% \item Skriv ämnen om programmering och systemutveckling som du är \Emph{nyfiken} på så tar jag upp dem i mån av tid. 
% \item Jag kommer också ge tips inför valfria tentan.  
% \item OBS! Inga föreläsningar w14.
% \end{itemize}

% \end{SlideExtra}


% % \Subsection{Grumligt-lådan \& Nyfiken-på-lådan}
% % \begin{Slide}{Grumligt-lådan och Nyfiken-på-lådan}
% % \begin{itemize}
% % \item Skriv lapp i \Alert{GRUMLIGT}-lådan om du har något \Alert{grundläggande begrepp} i kursen som du fortfarande tycker är \Alert{svårt att begripa}.
% %
% % \item[]
% %
% % \item Skriv lapp i \Emph{NYFIKEN-PÅ}-lådan om du vill veta mer om något ämne inom programmering och som går \Emph{bortom grunderna}.
% % \end{itemize}
% % \end{Slide}

% \fi